\documentclass[10pt]{article}

\usepackage[usenames]{color} %used for font color
\usepackage{amssymb} %maths
\usepackage{amsmath} %maths
\usepackage[utf8]{inputenc} %useful to type directly diacritic characters
\usepackage{physics}
\usepackage[capitalize]{cleveref} % Load last, especially AFTER algorithm2e.
\Crefname{figure}{Fig.}{Figs.}
\crefrangeformat{line}{lines~#3#1#4--#5#2#6}
\crefformat{equation}{(#2#1#3)}
\crefrangeformat{equation}{(#3#1#4--#5#2#6)}
\newcommand{\creflastconjunction}{, and\nobreakspace}

\newtheorem{theorem}{Theorem}[section]
\newtheorem{lemma}[theorem]{Lemma}

\begin{document}

\section{Proof for cascaded amplitude kickback for expectation comparison (1st-moment, or weighted mean)}

\begin{theorem}
In a quantum system with two hilbert spaces, one for actions $a$ and the other for retuns $z$, denoted by $\mathcal{H}_{\mathcal{A}} \otimes \mathcal{H}_{\mathcal{Z}} = \mathbb{C}^{2^{|\mathcal{A}|}} \otimes \mathbb{C}^{2^{|\mathcal{Z}|}}$ where $a \in \mathcal{A}$ is the set of actions and $z \in \mathcal{Z}$ is the set of return atoms, the $m$-th moment of the return atoms can be compared between actions entirely in quantum space, i.e., without needing to encode/decode between classical computation.
\end{theorem}

\begin{proof}

A comparison of the $m$-th moments can be done entirely in quantum space by first applying a cascade of return-atom preparations conditioned on the action space, followed by a reset over the return-atoms space, producing a weighted sum of $m$-th moment expectations over the return-atoms raised to the power of the number of cascades.

We start with an initial state of all zeros across the two spaces, which looks like 

\begin{equation}
\ket{\psi_{\textrm{in}}} = \ket{0}^{\otimes |\mathcal{A}|} \otimes \ket{0}^{\otimes |\mathcal{Z}|}
\end{equation}

We prepare an initial state of superposition for the actions via some operator $U_{a}$:

\begin{equation}
\ket{\psi_0} = (U_{a} \otimes \mathbb{I}) \ket{\psi_{\textrm{in}}} = \sum_{a} \sqrt{p(a)} \ket{a} \otimes \ket{0}^{\otimes |\mathcal{Z}|}.
\end{equation}

For example, if an equal superposition of actions is desired, i.e., $p(a) = 1/|\mathcal{A}|$, one could select $U_{a} = H^{\otimes |\mathcal{A}|}$, which is simply the Hadamard gate on every action qubit.

It will become convenient to represent this state in density matrix form, for reasons we will see shortly. This state can be represented as the density matrix:

\begin{align}
\label{eq:rho_0}
\rho_0
&= \ket{\psi_0}\bra{\psi_0} \\
&= (\sum_{a_i} \sqrt{p(a_i)} \ket{a_i} \otimes \ket{0}^{\otimes |\mathcal{Z}|}) (\sum_{a_j} \sqrt{p(a_j)} \bra{a_j} \otimes \bra{0}^{\otimes |\mathcal{Z}|}) \\
&= \sum_{a_i} \sum_{a_j} \sqrt{p(a_i)} \sqrt{p(a_j)} \ket{a_i} \otimes \ket{0}^{\otimes |\mathcal{Z}|} \bra{a_i} \otimes \bra{0}^{\otimes |\mathcal{Z}|}.
\end{align}

Next, we define an operator over the combined Hilbert space $\mathcal{H}_{\mathcal{A}} \otimes \mathcal{H}_{\mathcal{Z}}$ for preparing a probability distribution of return-atoms for every action -- this would be the ML model in a distributional reinforcement learning setting -- called $U_{z|a}$

\begin{align}
\label{eq:U_z_a}
U_{z|a}
    &= \sum_{a} \ket{a}\bra{a} \otimes U_{z}(a) \\
    &= \sum_{a} \ket{a}\bra{a} \otimes \Bigl[ \sum_{z} \sqrt{p(z|a)} \ket{z}\bra{0} \Bigr],
\end{align}

which, when operating on the state-vector form of the system state, behaves as follows

\begin{equation}
U_{z|a} \ket{\psi_0} = \sum_{a} \sqrt{p(a)} \ket{a} \otimes \Bigl[ \sum_{z} \sqrt{p(z|a)} \ket{z} \Bigr].
\end{equation}

and in density matrix form

\begin{align}
U_{z|a} \rho_0 U_{z|a}^{\dagger}
&= U_{z|a} \ket{\psi_0} \bra{\psi_0} U_{z|a}^{\dagger} \\
&= \Bigl( \sum_{a_i} \sqrt{p(a_i)} \ket{a_i} \otimes \Bigl[ \sum_{z_j} \sqrt{p(z_j|a_i)} \ket{z_j} \Bigr] \Bigr) \Bigl( \sum_{a_k} \sqrt{p(a_k)} \bra{a_k} \otimes \Bigl[ \sum_{z_l} \sqrt{p(z_l|a_k)} \bra{z_l} \Bigr] \Bigr) \\
&= \sum_{a_i} \sum_{z_j} \sum_{a_k} \sum_{z_l} \Bigl( \sqrt{p(a_i)} \sqrt{p(z_j|a_i)} \sqrt{p(a_k)} \sqrt{p(z_l|a_k)} \Bigr) \ket{a_i} \ket{z_j} \bra{a_k} \bra{z_l}.
\end{align}

Note that the definition we provide for $U_{z|a}$ is not strictly unitary because we only include the outer $\bra{0}$ term -- that is, the effect of applying $U_{z|a}$ on an all-zero input state in the $\mathcal{H}_{\mathcal{Z}}$ space. While a complete definition for $U_{z|a}$ will also include effects of all other non-zero states in the $\mathcal{H}_{\mathcal{Z}}$ space (i.e., $\{\bra{i}\}_{i=1}^{2^{|\mathcal{Z}|}-1}$) their effect, as we will see shortly, will be inconsequential as we only consider systems that have a known all-zero starting state. We therefore omit these extra terms in the derivations for brevity.

Now we define an operator over the return-atoms Hilbert space $\mathcal{H}_{\mathcal{Z}}$ that encodes the return-atom values called $U_{v|z}$

\begin{align}
\label{eq:U_v_z}
U_{v|z} 
    = \sum_{z} \sqrt{v_z} \ket{z}\bra{z}
    = \begin{bmatrix}
    \sqrt{v_0} & 0 & 0 & \dots \\
    0 & \sqrt{v_1} & 0 &  \\
    0 & 0 & \sqrt{v_2} &  \\
    \vdots &  &  & \ddots
    \end{bmatrix},
\end{align}

which is a diagonal matrix, where $\sum_{z} v_z = 1$, i.e., the return-atom values form a probability distribution. Because the return-atom values form a probability distribution, $U_{v|z}$ is both *completely positive* and $U_{v|z}^{\dagger} U_{v|z} \leq \mathbb{I}$, i.e., *trace non-increasing*, which means that $U_{v|z}$ is a valid Kraus operator and can therefore be implemented on quantum hardware.

Notice that $U_{v|z}$ is Hermitian (i.e., eigenvalues are real), but not unitary. This means that the resulting state after applying $U_{v|z}$ will not be normalized. This normalization, however, can be circumvented by applying a reset operation on the return-atom space, as we will see shortly.

Applying $m$ repetitions of $U_{v|z}$ on $U_{z|a} \ket{\psi_0}$ produces the state

\begin{align}
\Bigl[ \prod_{m} (\mathbb{I} \otimes U_{v|z}) \Bigr]  U_{z|a} \ket{\psi_0} 
    &= \sum_{a} \sqrt{p(a)} \ket{a} \otimes \Bigl[ \sum_{z} \sqrt{p(z|a)} (\sqrt{v_z})^{m} \ket{z} \Bigr]
\end{align}

and again in density matrix form

\begin{align}
&\Bigl[ \prod_{m} (\mathbb{I} \otimes U_{v|z}) \Bigr]  U_{z|a} \rho_0 U_{z|a}^{\dagger} {\Bigl[ \prod_{m} (\mathbb{I} \otimes U_{v|z}) \Bigr]}^{\dagger}
\\&= \Bigl[ \prod_{m} (\mathbb{I} \otimes U_{v|z}) \Bigr]  U_{z|a} \ket{\psi_0} \bra{\psi_0} U_{z|a}^{\dagger} {\Bigl[ \prod_{m} (\mathbb{I} \otimes U_{v|z}) \Bigr]}^{\dagger}
\\&= \Bigl( \sum_{a_i} \sqrt{p(a_i)} \ket{a_i} \otimes \Bigl[ \sum_{z_j} \sqrt{p(z_j|a_i)} (\sqrt{v_{z_j}})^{m} \ket{z_j} \Bigr] \Bigr) \Bigl( \sum_{a_k} \sqrt{p(a_k)} \bra{a_k} \otimes \Bigl[ \sum_{z_l} \sqrt{p(z_l|a_k)} (\sqrt{v_{z_l}})^{m} \bra{z_l} \Bigr] \Bigr)
\\&= \sum_{a_i} \sum_{z_j} \sum_{a_k} \sum_{z_l} \Bigl( \sqrt{p(a_i)} \sqrt{p(z_j|a_i)} (\sqrt{v_{z_j}})^{m} \sqrt{p(a_k)} \sqrt{p(z_l|a_k)} (\sqrt{v_{z_l}})^{m} \Bigr) \ket{a_i} \ket{z_j} \bra{a_k} \bra{z_l}
\\&= \rho_1 \label{eq:rho_1}.
\end{align}

which we will refer to as $\rho_1$, where the subscript denotes the number of times the operations $ \Bigl[ \prod_{m} (\mathbb{I} \otimes U_{v|z}) \Bigr]  U_{z|a}$ have been applied -- this subscript index will be important later.

Now, to correct the normalization on the return-atoms space, we apply a reset operation over the space of the return-atoms $\mathcal{H}_{\mathcal{Z}}$ to collapse their qubits. A reset operation projects the selection of qubits to the $\ket{0}$ state. To accomplish this we apply a mid-circuit measurement on the return-atom qubits ($\mathcal{H}_{\mathcal{Z}}$) to collapse their state to to one of $\ket{z} \in \{\ket{z}\}_{z\in [1,|\mathcal{Z}|]}$ and then apply a conditional correction to map $\ket{z} \to \ket{0}$ (e.g., a bit-flip operation if $z \in \{0,1\}$). Mathematically, we can represent this collapse as a measurement projection followed by a partial-trace over the $\mathcal{H}_{\mathcal{Z}}$ space. Importantly, in this setting we do not care about the measurement result, only the collapse is desired.

Before we apply the reset to our system state, we first need to define the measurement process. A quantum measurement $M$ maps a quantum state in density matrix form $\rho$ to a mixed state using a set of projection operators $\Pi$ outlined by the measurement basis, given by

\begin{equation}
\rho' = M_{\mathcal{B}}[\rho] = \sum_{b} \Pi_b \rho \Pi^{\dagger}_b
\end{equation}

where $\rho$ is the density matrix for a quantum system of $q$ qubits, and $b \in \mathcal{B}$ is the measurement basis -- for example, if measuring in the Pauli-$Z$ basis, $\Pi_b = \ket{b}\bra{b}$, $\forall b \in \{0,2^{q}-1\}$. This resulting density matrix after measurement must also be normalized, which we can do by dividing by its trace (which is just the sum of its diagonal elements) as given by

\begin{equation}
\hat{\rho}' = \frac{\rho'}{\textrm{tr}[\rho']}.
\end{equation}

Since $\hat{\rho}'$ spans the space of both actions and return-atoms $\mathcal{H}_{\mathcal{A}} \otimes \mathcal{H}_{\mathcal{Z}}$, to complete the reset over the return-atoms space after the projective measurement we must trace-out the return-atoms qubits and then simply set their state to zero. To do this, we define the completed reset operation

\begin{align}
\textrm{Reset}_{\mathcal{Z}} [\rho] 
    &= \textrm{tr}_{\mathcal{Z}}{\bigl[ \hat{\rho}' \bigr]} \otimes \ket{0}^{\otimes |\mathcal{Z}|}\bra{0}^{\otimes |\mathcal{Z}|} \\
    &= \textrm{tr}_{\mathcal{Z}}{\Bigl[ \frac{M_{\mathcal{Z}}[\rho]}{\textrm{tr}\bigl[M_{\mathcal{Z}}[\rho]\bigr]} \Bigr]} \otimes \ket{0}^{\otimes |\mathcal{Z}|}\bra{0}^{\otimes |\mathcal{Z}|} \\
    &= \tilde{\rho}
\end{align}

where $\textrm{tr}_{\mathcal{Z}} : \mathcal{H}_{\mathcal{A}} \otimes \mathcal{H}_{\mathcal{Z}} \to \mathcal{H}_{\mathcal{A}}$ is the partial trace operator which traces out (i.e., removes) the subspace $\mathcal{H}_{\mathcal{Z}}$.

So, applying this to our distributional RL setting, we can resume from $\rho_1$ and reset the return-atom qubits in the following manner. First, we define the set of projectors $\Pi = \{ \mathbb{I}\otimes\ket{z}\bra{z} \}_{z\in \mathcal{Z}}$, which covers all possible return-atom outcomes while doing nothing to the action space, and then apply these projectors as a *partial* measurement on $\rho_1$ as follows

\begin{align}
M_{\mathcal{Z}}[\rho_1]
    &= \sum_{z} \Pi_z \rho_1 \Pi^{\dagger}_z
    \\&= \sum_{z} \bigl(\mathbb{I}\otimes\ket{z}\bra{z}\bigr) \rho_1 \bigl(\mathbb{I}\otimes\ket{z}\bra{z}\bigr)
    \\&= \sum_{z} \sum_{a_i} \sum_{a_k} \Bigl( \sqrt{p(a_i)} \sqrt{p(z|a_i)} (\sqrt{v_{z}})^{m} \sqrt{p(a_k)} \sqrt{p(z|a_k)} (\sqrt{v_{z}})^{m} \Bigr) \ket{a_i} \ket{z} \bra{a_k} \bra{z}
    \\&~~~~~\textrm{(Collapses where} ~ z_j=z_l=z \textrm{)}
    \\&= \rho'_1
\end{align}

which has the normalization factor

\begin{align}
\textrm{tr}\bigl[\rho'_1\bigr]
    &= \sum_{z} \sum_{a} \sqrt{p(a)} \sqrt{p(z|a)} (\sqrt{v_{z}})^{m} \sqrt{p(a)} \sqrt{p(z|a)} (\sqrt{v_{z}})^{m}
    \\&~~~~~\textrm{(Sum along diagonal collapses where} ~ a_i=a_k=a \textrm{)}
    \\&= \sum_{z} \sum_{a} p(a) p(z|a) v_{z}^{m}
    \\&= \sum_{a} p(a) \sum_{z} p(z|a) v_{z}^{m}
    \\&= \sum_{a} p(a) \mathbb{E}[Z^{m} | a]
    \\&~~~~~\textrm{(By definition of expectation of a random variable)}.
\end{align}

Notice that here is where the expectation of the random return conditioned on the action space, $\mathbb{E}[Z^{m} | a]$, first appears -- importantly, we have the $m$-th moment of $Z$ conditioned on $a$. Thus, the normalized density matrix is

\begin{equation}
{\hat{\rho}'}_1
= \frac{\rho'_1}{\textrm{tr}\bigl[\rho'_1\bigr]}
= (\frac{1}{\sum_{a} p(a) \mathbb{E}[Z^{m} | a]}) \cdot \sum_{z} \sum_{a_i} \sum_{a_k} \Bigl( \sqrt{p(a_i)} \sqrt{p(z|a_i)} (\sqrt{v_{z}})^{m} \sqrt{p(a_k)} \sqrt{p(z|a_k)} (\sqrt{v_{z}})^{m} \Bigr) \ket{a_i} \ket{z} \bra{a_k} \bra{z}
\end{equation}

Next, we can find the reduced state of the action space after the partial measurement by taking the partial trace over the return-atoms space

\begin{align}
\textrm{tr}_{\mathcal{Z}}{\bigl[ {\hat{\rho}'}_1 \bigr]}
&= \sum_{z} (\mathbb{I}\otimes\bra{z}) {\hat{\rho}'}_1 (\mathbb{I}\otimes\ket{z}) \\
&= (\frac{1}{\sum_{a} p(a) \mathbb{E}[Z^{m} | a]}) \cdot \sum_{z} \sum_{a_i} \sum_{a_k} \Bigl( \sqrt{p(a_i)} \sqrt{p(z|a_i)} (\sqrt{v_{z}})^{m} \sqrt{p(a_k)} \sqrt{p(z|a_k)} (\sqrt{v_{z}})^{m} \Bigr) \ket{a_i}\bra{a_k} .
\end{align}

The final system state after applying the reset operation is therefore

\begin{align}
\textrm{Reset}_{\mathcal{Z}} [\rho_1] 
    &= \textrm{tr}_{\mathcal{Z}}{\bigl[ {\hat{\rho}'}_1 \bigr]} \otimes \ket{0}^{\otimes |\mathcal{Z}|}\bra{0}^{\otimes |\mathcal{Z}|} \\
    &= (\frac{1}{\sum_{a} p(a) \mathbb{E}[Z^{m} | a]}) \cdot \sum_{z} \sum_{a_i} \sum_{a_k} \Bigl( \sqrt{p(a_i)} \sqrt{p(z|a_i)} (\sqrt{v_{z}})^{m} \sqrt{p(a_k)} \sqrt{p(z|a_k)} (\sqrt{v_{z}})^{m} \Bigr) \ket{a_i}\bra{a_k} \otimes \ket{0}^{\otimes |\mathcal{Z}|}\bra{0}^{\otimes |\mathcal{Z}|} \\
    &= \tilde{\rho_1}
\end{align}

Notice that at this point the return-atoms space has been reduced to the all-zero state $\ket{0}^{\otimes |\mathcal{Z}|}$, which has a very similar form to $\rho_0$ in \cref{eq:rho_0}. This means that, after applying the reset, the system state is placed back into a form after applying superposition of actions, and is thus receptive to the sequential return-atom operations $(\mathbb{I} \otimes U_{v|z}) U_{z|a}$.

If at this point we were to stop and measure the action qubits using projectors $\Pi = \{\ket{a}\bra{a}\}_{a \in \mathcal{A}}$ -- ignoring the return-atom space because it is already zero -- we would get the following state in density matrix form

\begin{align}
M_{\mathcal{A}}[\tilde{\rho_1}]
&= M_{\mathcal{A}}[\textrm{tr}_{\mathcal{Z}}{\bigl[ {\hat{\rho}'}_1 \bigr]}] 
\\&~~~~~\textrm{(Note that inclusion of partial trace is for formal removal of return-atom space)}
\\&= \sum_{a} \Pi_a \textrm{tr}_{\mathcal{Z}}{\bigl[ {\hat{\rho}'}_1 \bigr]} \Pi^{\dagger}_a
\\&= \sum_{a} \ket{a}\bra{a} \textrm{tr}_{\mathcal{Z}}{\bigl[ {\hat{\rho}'}_1 \bigr]} \ket{a}\bra{a}
\\&= \sum_{a} \ket{a}\bra{a} \Bigl((\frac{1}{\sum_{a} p(a) \mathbb{E}[Z^{m} | a]}) \cdot \sum_{z} \sum_{a_i} \sum_{a_k} \bigl( \sqrt{p(a_i)} \sqrt{p(z|a_i)} (\sqrt{v_{z}})^{m} \sqrt{p(a_k)} \sqrt{p(z|a_k)} (\sqrt{v_{z}})^{m} \bigr) \ket{a_i}\bra{a_k} \Bigr) \ket{a}\bra{a}
\\&= (\frac{1}{\sum_{a} p(a) \mathbb{E}[Z^{m} | a]}) \cdot \sum_{a} \sum_{z} \bigl( \sqrt{p(a_i=a)} \sqrt{p(z|a_i=a)} (\sqrt{v_{z}})^{m} \sqrt{p(a_k=a)} \sqrt{p(z|a_k=a)} (\sqrt{v_{z}})^{m} \bigr) \ket{a}\bra{a} 
\\&~~~~~\textrm{(Sums collapse when} ~ a_i = a_k = a \textrm{)}
\\&= (\frac{1}{\sum_{a} p(a) \mathbb{E}[Z^{m} | a]}) \cdot \sum_{a} \sum_{z} p(a) p(z|a) v_{z}^{m} \ket{a}\bra{a}
\\&= (\frac{1}{\sum_{a} p(a) \mathbb{E}[Z^{m} | a]}) \cdot \sum_{a} p(a) \sum_{z} p(z|a) v_{z}^{m} \ket{a}\bra{a}
\\&= (\frac{1}{\sum_{a} p(a) \mathbb{E}[Z^{m} | a]}) \cdot \sum_{a} p(a) \mathbb{E}[Z^m | a] \ket{a}\bra{a} 
\\&~~~~~\textrm{(By definition of expectation of a random variable)},
\end{align}

which is diagonal with normalization factor

\begin{align}
\textrm{tr}\Bigl[ M_{\mathcal{A}}[\textrm{tr}_{\mathcal{Z}}{\bigl[ {\hat{\rho}'}_1 \bigr]}] \Bigr] = (\frac{1}{\sum_{a} p(a) \mathbb{E}[Z^{m} | a]}) \cdot \sum_{a} p(a) \mathbb{E}[Z^{m} | a] ~~\textrm{(Sum of diagonal elements)},
\end{align}

producing the normalized state

\begin{align}
\frac{M_{\mathcal{A}}[\textrm{tr}_{\mathcal{Z}}{\bigl[ {\hat{\rho}'}_1 \bigr]}]}{\textrm{tr}\Bigl[ M_{\mathcal{A}}[\textrm{tr}_{\mathcal{Z}}{\bigl[ {\hat{\rho}'}_1 \bigr]}] \Bigr]}
&= (\frac{1}{(\frac{1}{\sum_{a} p(a) \mathbb{E}[Z^{m} | a]}) \cdot \sum_{a} p(a) \mathbb{E}[Z^{m} | a]}) \cdot (\frac{1}{\sum_{a} p(a) \mathbb{E}[Z^{m} | a]}) \cdot \sum_{a} p(a) \mathbb{E}[Z^{m} | a] \ket{a}\bra{a} \\
&= (\frac{1}{\sum_{a} p(a) \mathbb{E}[Z^{m} | a]}) \cdot \sum_{a} p(a) \mathbb{E}[Z^{m} | a] \ket{a}\bra{a}.
\end{align}

Notice that because the state was normalized when creating ${\hat{\rho}'}_1$ this state when measuring $a$ was already normalized, but for completeness we complete the process again here.

This means that the probability of measuring an arbitrary action $a_i$ is

\begin{equation}
P[\ket{a_i}] = \frac{p(a_i) \mathbb{E}[Z^{m} | a_i]}{\sum_{a} p(a) \mathbb{E}[Z^{m} | a]},
\end{equation}

which is the \emph{normalized weighted expectation of return-atoms for the given action $a_i$ relative to all other actions $\mathcal{A} \setminus \{a_i\}$}.



% Iteration 2

If we do not measure the action space, and instead apply another iteration of cascaded return-atom operations $(\mathbb{I} \otimes U_{v|z}) U_{z|a}$, we get the following state in density matrix form

\begin{align}
& \Bigl[ \prod_{m} (\mathbb{I} \otimes U_{v|z}) \Bigr] U_{z|a} \tilde{\rho_1} U_{z|a}^{\dagger} {\Bigl[ \prod_{m} (\mathbb{I} \otimes U_{v|z}) \Bigr]}^{\dagger}
\\&= \Bigl[ \prod_{m} (\mathbb{I} \otimes U_{v|z}) \Bigr] U_{z|a} \textrm{Reset}_{\mathcal{Z}} [\rho_1] U_{z|a}^{\dagger} {\Bigl[ \prod_{m} (\mathbb{I} \otimes U_{v|z}) \Bigr]}^{\dagger}
% 
\\&= \Bigl[ \prod_{m} (\mathbb{I} \otimes U_{v|z}) \Bigr] U_{z|a} \Bigl[ (\frac{1}{\sum_{a} p(a) \mathbb{E}[Z^{m} | a]}) \cdot \sum_{z} \sum_{a_i} \sum_{a_k} \Bigl( \sqrt{p(a_i)} \sqrt{p(z|a_i)} (\sqrt{v_{z}})^{m} \sqrt{p(a_k)} \sqrt{p(z|a_k)} (\sqrt{v_{z}})^{m} \Bigr) \ket{a_i}\bra{a_k} \otimes \ket{0}^{\otimes |\mathcal{Z}|}\bra{0}^{\otimes |\mathcal{Z}|} \Bigr] U_{z|a}^{\dagger} {\Bigl[ \prod_{m} (\mathbb{I} \otimes U_{v|z}) \Bigr]}^{\dagger}
% 
\\&= (\frac{1}{\sum_{a} p(a) \mathbb{E}[Z^{m} | a]}) \cdot \Bigl[ \prod_{m} (\mathbb{I} \otimes U_{v|z}) \Bigr] U_{z|a} \Bigl[ \sum_{z} \sum_{a_i} \sum_{a_k} \Bigl( \sqrt{p(a_i)} \sqrt{p(z|a_i)} (\sqrt{v_{z}})^{m} \sqrt{p(a_k)} \sqrt{p(z|a_k)} (\sqrt{v_{z}})^{m} \Bigr) \ket{a_i}\bra{a_k} \otimes \ket{0}^{\otimes |\mathcal{Z}|}\bra{0}^{\otimes |\mathcal{Z}|} \Bigr] U_{z|a}^{\dagger} {\Bigl[ \prod_{m} (\mathbb{I} \otimes U_{v|z}) \Bigr]}^{\dagger}
\\&~~~~~\textrm{(Pull out constant normalization factor)}
% 
\\&= (\frac{1}{\sum_{a} p(a) \mathbb{E}[Z^{m} | a]}) \cdot \Bigl[ \prod_{m} (\mathbb{I} \otimes U_{v|z}) \Bigr] U_{z|a} \Bigl[ \sum_{z_q} \sum_{a_i} \sum_{a_k} \Bigl( \sqrt{p(a_i)} \sqrt{p(z_q|a_i)} (\sqrt{v_{z_q}})^{m} \sqrt{p(a_k)} \sqrt{p(z_q|a_k)} (\sqrt{v_{z_q}})^{m} \Bigr) \ket{a_i}\bra{a_k} \otimes \ket{0}^{\otimes |\mathcal{Z}|}\bra{0}^{\otimes |\mathcal{Z}|} \Bigr] U_{z|a}^{\dagger} {\Bigl[ \prod_{m} (\mathbb{I} \otimes U_{v|z}) \Bigr]}^{\dagger}
\\&~~~~~\textrm{(Simple relabeling the outer-sum return-atom index}~ z \to z_q \textrm{)}
% 
\\&= (\frac{1}{\sum_{a} p(a) \mathbb{E}[Z^{m} | a]}) \cdot \Bigl[ \sum_{z_q} \sum_{a_i} \sum_{a_k} \sum_{z_w} \Bigl( \sqrt{p(a_i)} \sqrt{p(z_q|a_i)} (\sqrt{v_{z_q}})^{m} \sqrt{p(a_k)} \sqrt{p(z_q|a_k)} (\sqrt{v_{z_q}})^{m} \sqrt{p(z_w|a_i)} (\sqrt{v_{z_w}})^{m} \Bigr) \ket{a_i}\bra{a_k} \otimes \ket{z_w}\bra{0}^{\otimes |\mathcal{Z}|} \Bigr] (\mathbb{I} \otimes U_{v|z_q}^{\dagger})
\\&~~~~~\textrm{(LHS operators collapse where}~ a_i = a ~\textrm{from}~ U_{z|a} ~\textrm{and create new index where}~ z_{w} = z ~\textrm{from}~ U_{v|z} ~\textrm{)}
% 
\\&= (\frac{1}{\sum_{a} p(a) \mathbb{E}[Z^{m} | a]}) \cdot \Bigl[ \sum_{z_q} \sum_{a_i} \sum_{a_k} \sum_{z_w} \sum_{z_t} \Bigl( \sqrt{p(a_i)} \sqrt{p(z_q|a_i)} (\sqrt{v_{z_q}})^{m} \sqrt{p(a_k)} \sqrt{p(z_q|a_k)} (\sqrt{v_{z_q}})^{m} \sqrt{p(z_w|a_i)} (\sqrt{v_{z_w}})^{m} \sqrt{p(z_t|a_k)} (\sqrt{v_{z_t}})^{m}  \Bigr) \ket{a_i}\bra{a_k} \otimes \ket{z_w}\bra{z_t} \Bigr]
\\&~~~~~\textrm{(RHS operators collapse where}~ a_k = a ~\textrm{from}~ U_{z|a}^{\dagger} ~\textrm{and create new index where}~ z_{t} = z ~\textrm{from}~ U_{v|z}^{\dagger} ~\textrm{)}
% 
\\&= \rho_2 \label{eq:rho_2}.
\end{align}

We then apply the reset operation to the return-atoms space by first applying return-atoms projectors 

\begin{align}
M_{\mathcal{Z}}[\rho_2]
    &= \sum_{z} \Pi_z \rho_2 \Pi^{\dagger}_z \\
    &= \sum_{z} \bigl(\mathbb{I}\otimes\ket{z}\bra{z}\bigr) \rho_2 \bigl(\mathbb{I}\otimes\ket{z}\bra{z}\bigr) \\
    &= (\frac{1}{\sum_{a} p(a) \mathbb{E}[Z^{m} | a]}) \cdot \Bigl[ \sum_{z} \sum_{z_q} \sum_{a_i} \sum_{a_k} \Bigl( \sqrt{p(a_i)} \sqrt{p(z_q|a_i)} (\sqrt{v_{z_q}})^{m} \sqrt{p(a_k)} \sqrt{p(z_q|a_k)} (\sqrt{v_{z_q}})^{m} \sqrt{p(z|a_i)} (\sqrt{v_{z}})^{m} \sqrt{p(z|a_k)} (\sqrt{v_{z}})^{m}  \Bigr) \ket{a_i}\bra{a_k} \otimes \ket{z}\bra{z} \Bigr]
    \\&~~~~~\textrm{(Collapses where}~ z_w = z_t = z ~\textrm{)} \\
    &= \rho'_2 
\end{align}

which has the normalization factor

\begin{align}
\textrm{tr}\bigl[\rho'_2\bigr]
    &= (\frac{1}{\sum_{a} p(a) \mathbb{E}[Z^{m} | a]}) \cdot \Bigl[ \sum_{z} \sum_{z_q} \sum_{a} \sqrt{p(a)} \sqrt{p(z_q|a)} (\sqrt{v_{z_q}})^{m} \sqrt{p(a)} \sqrt{p(z_q|a)} (\sqrt{v_{z_q}})^{m} \sqrt{p(z|a)} (\sqrt{v_{z}})^{m} \sqrt{p(z|a)} (\sqrt{v_{z}})^{m} \Bigr]
    \\&~~~~~\textrm{(Sum along diagonal collapses where}~ a_i = a_k = a ~\textrm{)}
    \\&= (\frac{1}{\sum_{a} p(a) \mathbb{E}[Z^{m} | a]}) \cdot \Bigl[ \sum_{z} \sum_{z_q} \sum_{a} p(a) p(z_q|a) v_{z_q}^{m} p(z|a) v_{z}^{m} \Bigr]
    \\&= (\frac{1}{\sum_{a} p(a) \mathbb{E}[Z^{m} | a]}) \cdot \Bigl[ \sum_{a} p(a) \sum_{z} p(z|a) v_{z}^{m} \sum_{z_q} p(z_q|a) v_{z_q}^{m} \Bigr]
    \\&= (\frac{1}{\sum_{a} p(a) \mathbb{E}[Z^{m} | a]}) \cdot \Bigl[ \sum_{a} p(a) \bigl( \sum_{z} p(z|a) v_{z}^{m} \bigr)^2 \Bigr]
    \\&~~~~~\textrm{(Both return-atom sums are identical, i.e.,}~ z = z_q ~\textrm{)}
    \\&= (\frac{1}{\sum_{a} p(a) \mathbb{E}[Z^{m} | a]}) \cdot \Bigl[ \sum_{a} p(a) \mathbb{E}[Z^{m} | a]^2 \Bigr]
    \\&~~~~~\textrm{(By definition of expectation of a random variable)}.
\end{align}

Thus the normalized density matrix is

\begin{align}
{\hat{\rho}'}_2
    &= \frac{\rho'_2}{\textrm{tr}\bigl[\rho'_2\bigr]}
    \\&= \frac{(\frac{1}{\sum_{a} p(a) \mathbb{E}[Z^{m} | a]}) \cdot \Bigl[ \sum_{z} \sum_{z_q} \sum_{a_i} \sum_{a_k} \Bigl( \sqrt{p(a_i)} \sqrt{p(z_q|a_i)} (\sqrt{v_{z_q}})^{m} \sqrt{p(a_k)} \sqrt{p(z_q|a_k)} (\sqrt{v_{z_q}})^{m} \sqrt{p(z|a_i)} (\sqrt{v_{z}})^{m} \sqrt{p(z|a_k)} (\sqrt{v_{z}})^{m}  \Bigr) \ket{a_i}\bra{a_k} \otimes \ket{z}\bra{z} \Bigr]}{(\frac{1}{\sum_{a} p(a) \mathbb{E}[Z^{m} | a]}) \cdot \Bigl[ \sum_{a} p(a) \mathbb{E}[Z^{m} | a]^2 \Bigr]}
    \\&= (\frac{1}{\sum_{a} p(a) \mathbb{E}[Z^{m} | a]^2}) \cdot \Bigl[ \sum_{z} \sum_{z_q} \sum_{a_i} \sum_{a_k} \Bigl( \sqrt{p(a_i)} \sqrt{p(z_q|a_i)} (\sqrt{v_{z_q}})^{m} \sqrt{p(a_k)} \sqrt{p(z_q|a_k)} (\sqrt{v_{z_q}})^{m} \sqrt{p(z|a_i)} (\sqrt{v_{z}})^{m} \sqrt{p(z|a_k)} (\sqrt{v_{z}})^{m}  \Bigr) \ket{a_i}\bra{a_k} \otimes \ket{z}\bra{z} \Bigr].
\end{align}

Tracing out the return-atoms system gives

\begin{align}
\textrm{tr}_{\mathcal{Z}}{\bigl[ {\hat{\rho}'}_2 \bigr]}
    &= \sum_{z} (\mathbb{I}\otimes\bra{z}) {\hat{\rho}'}_2 (\mathbb{I}\otimes\ket{z})
    \\&= (\frac{1}{\sum_{a} p(a) \mathbb{E}[Z^{m} | a]^2}) \cdot \Bigl[ \sum_{z} \sum_{z_q} \sum_{a_i} \sum_{a_k} \Bigl( \sqrt{p(a_i)} \sqrt{p(z_q|a_i)} (\sqrt{v_{z_q}})^{m} \sqrt{p(a_k)} \sqrt{p(z_q|a_k)} (\sqrt{v_{z_q}})^{m} \sqrt{p(z|a_i)} (\sqrt{v_{z}})^{m} \sqrt{p(z|a_k)} (\sqrt{v_{z}})^{m}  \Bigr) \ket{a_i}\bra{a_k} \Bigr]
\end{align}

and the final system state after applying the reset operation is therefore

\begin{align}
\textrm{Reset}_{\mathcal{Z}} [\rho_2]
    &= \textrm{tr}_{\mathcal{Z}}{\bigl[ {\hat{\rho}'}_2 \bigr]} \otimes \ket{0}^{\otimes |\mathcal{Z}|}\bra{0}^{\otimes |\mathcal{Z}|}
    \\&= (\frac{1}{\sum_{a} p(a) \mathbb{E}[Z^{m} | a]^2}) \cdot \Bigl[ \sum_{z} \sum_{z_q} \sum_{a_i} \sum_{a_k} \Bigl( \sqrt{p(a_i)} \sqrt{p(z_q|a_i)} (\sqrt{v_{z_q}})^{m} \sqrt{p(a_k)} \sqrt{p(z_q|a_k)} (\sqrt{v_{z_q}})^{m} \sqrt{p(z|a_i)} (\sqrt{v_{z}})^{m} \sqrt{p(z|a_k)} (\sqrt{v_{z}})^{m}  \Bigr) \ket{a_i}\bra{a_k} \otimes \ket{0}^{\otimes |\mathcal{Z}|}\bra{0}^{\otimes |\mathcal{Z}|} \Bigr]
    \\&= \tilde{\rho_2}
\end{align}

Importantly, notice that a squared version of the expectation over return-atoms appears in the normalization term. From this we readily see a correlation between the power index of the expectation and the number of times the return-atoms operations have been applied (conveniently denoted in the subscript of the system state ${\hat{\rho}'}_2$). This is even more prevalent if we apply a measurement on the action space as follows

\begin{align}
M_{\mathcal{A}}[\tilde{\rho_2}]
    &= M_{\mathcal{A}}[\textrm{tr}_{\mathcal{Z}}{\bigl[ {\hat{\rho}'}_2 \bigr]}]
    \\&= \sum_{a} \Pi_a \textrm{tr}_{\mathcal{Z}}{\bigl[ {\hat{\rho}'}_2 \bigr]} \Pi^{\dagger}_a
    \\&= \sum_{a} \ket{a}\bra{a} \textrm{tr}_{\mathcal{Z}}{\bigl[ {\hat{\rho}'}_2 \bigr]} \ket{a}\bra{a}
    \\&= \sum_{a} \ket{a}\bra{a} \Bigr( (\frac{1}{\sum_{a} p(a) \mathbb{E}[Z^{m} | a]^2}) \cdot \Bigl[ \sum_{z} \sum_{z_q} \sum_{a_i} \sum_{a_k} \Bigl( \sqrt{p(a_i)} \sqrt{p(z_q|a_i)} (\sqrt{v_{z_q}})^{m} \sqrt{p(a_k)} \sqrt{p(z_q|a_k)} (\sqrt{v_{z_q}})^{m} \sqrt{p(z|a_i)} (\sqrt{v_{z}})^{m} \sqrt{p(z|a_k)} (\sqrt{v_{z}})^{m}  \Bigr) \ket{a_i}\bra{a_k} \Bigr] \Bigl) \ket{a}\bra{a}
    \\&= (\frac{1}{\sum_{a} p(a) \mathbb{E}[Z^{m} | a]^2}) \cdot \Bigl[ \sum_{z} \sum_{z_q} \sum_{a} \Bigl( \sqrt{p(a)} \sqrt{p(z_q|a)} (\sqrt{v_{z_q}})^{m} \sqrt{p(a)} \sqrt{p(z_q|a)} (\sqrt{v_{z_q}})^{m} \sqrt{p(z|a)} (\sqrt{v_{z}})^{m} \sqrt{p(z|a)} (\sqrt{v_{z}})^{m}  \Bigr) \ket{a}\bra{a} \Bigr]
    \\&~~~~~\textrm{(Sums collapse when} ~ a_i = a_k = a \textrm{)}
    \\&= (\frac{1}{\sum_{a} p(a) \mathbb{E}[Z^{m} | a]^2}) \cdot \Bigl[ \sum_{z} \sum_{z_q} \sum_{a} p(a) p(z_q|a) v_{z_q}^{m} p(z|a) v_{z}^{m} \ket{a}\bra{a} \Bigr]
    \\&= (\frac{1}{\sum_{a} p(a) \mathbb{E}[Z^{m} | a]^2}) \cdot \Bigl[ \sum_{a} p(a) \sum_{z} p(z|a) v_{z}^{m} \sum_{z_q} p(z_q|a) v_{z_q}^{m} \ket{a}\bra{a}\Bigr]
    \\&= (\frac{1}{\sum_{a} p(a) \mathbb{E}[Z^{m} | a]^2}) \cdot \Bigl[ \sum_{a} p(a) \bigl( \sum_{z} p(z|a) v_{z}^{m} \bigr)^2 \ket{a}\bra{a} \Bigr]
    \\&~~~~~\textrm{(Both return-atom sums are identical, i.e.,}~ z = z_q ~\textrm{)}
    \\&= (\frac{1}{\sum_{a} p(a) \mathbb{E}[Z^{m} | a]^2}) \cdot \sum_{a} p(a) \mathbb{E}[Z^{m} | a]^2 \ket{a}\bra{a}
    \\&~~~~~\textrm{(By definition of expectation of a random variable)}.
\end{align}

Which we can readily see is both diagonal and normalized. The probability of measuring any action $a_i$ is therefore

\begin{equation}
P[\ket{a_i}] = \frac{p(a_i) \mathbb{E}[Z^{m} | a_i]^2}{\sum_{a} p(a) \mathbb{E}[Z^{m} | a]^2},
\end{equation}

which is the normalized weighted squared expectation of return atoms for the chosen action $a_i$ relative to all other actions.

Hence, there is a direct relationship between the number of times $\Bigl[ \prod_{m} (\mathbb{I} \otimes U_{v|z}) \Bigr] U_{z|a}$ has been applied and the power of the return-atom $m$-th moment expectation. Formally, we define the cascaded quantum measurement operation over the joint action-return space $\mathcal{H}_{\mathcal{A}} \otimes \mathcal{H}_{\mathcal{Z}}$ as

\begin{equation}
    M_{a} \Bigl( \prod_{n} \textrm{Reset}_{z} \Bigl[ \prod_{m} (\mathbb{I} \otimes U_{v|z}) \Bigr] U_{z|a} \Bigr) (U_{a} \otimes \mathbb{I}) : \ket{0}^{\otimes |\mathcal{A}|} \otimes \ket{0}^{\otimes |\mathcal{Z}|} \to (\frac{1}{\sum_{a} p(a) \mathbb{E}[Z^{m} | a]^{n}}) \cdot \sum_{a} p(a) \mathbb{E}[Z^{m} | a]^{n} \ket{a}\bra{a}, \label{eq:result}
\end{equation}

which pushes an initial all-zero state to a superposition of action states, where the probability of measuring any given action is the 
\emph{normalized weighted $m$-th moment expectation of return-atoms raised to the power of $n$, for any given action $a_i$ relative to all other actions $\mathcal{A} \setminus \{a_i\}$}.

\end{proof}

This result is important because it allows for comparison of the moments of many random variables within quantum space -- that is, the entire operation is completely quantum, without the need for encoding/decoding into the classical space mid-way through the algorithm, which is both expensive and error prone due to stochastic noise. The probability of measuring any given state in the action space is directly proportional to the $m$-th moment expectation of that action relative to all others. Moreover, we can boost the likelihood of measuring actions with higher expectations by intelligently selecting the cascade index $n$ to ``spread'' the expectation distribution, i.e., the probability higher expectations get raised while others get lowered. This allows interplay between the relative strength of action probabilities and the number of shots of the quantum circuit -- that is, higher values of $n$ can be chosen to lower the number of repeated shots required, and vice versa.

\section{Example Scenarios}

For example, consider a scenario where finding the action corresponding to the maximum average (mean) value of the return-atom distribution is desired, i.e., ${\arg\max}_{a}{\{\mathbb{E}[Z | a]\}}$. In this setting, we do not care about the actual expected return value, only the relative distance between expected returns. Applying \cref{eq:result} with $n=1$ and $m=1$ will yield a distribution across every 1st moment, i.e., the mean, of return-atoms. If the expectations across the actions are similar in value, then the likelihood of measuring any given action will be similar -- with a maximum value near $1/|\mathcal{A}|$. As we push $n \to \infty$, the relative (normalized) distance between the expectations becomes greater -- with a maximum value near $1$. Hence, the action associated with the maximum expected return can be found by choosing a sufficiently high $n$ such that the probability of measuring that action is boosted (approaching $1$) while others are lowered (approaching $0$) by only running a single shot of the quantum circuit.

Consider another example where instead we are interested in the similarity in expected returns, i.e., which actions have returns that are close to each other in relative magnitude. This could be beneficial in situations where exploration in the environment is desired, in contrast to a greedy policy (i.e., always takes the perceived ``best'' action). In this setting, one could choose a very low value for $n$ and instead run many shots of the circuit to generate a probability distribution over the action space, i.e., the set $\{ \frac{p(a_i)\mathbb{E}[Z | a_i]}{\sum_{a \in \mathcal{A}} p(a)\mathbb{E}[Z | a]} \}_{a_i \in \mathcal{A}}$.

As another example, consider the comparison of the 2nd moments of the returns distribution. Selecting $m=2$ with low $n$ one can sample the circuit many times to generate the distribution $\{ \frac{p(a_i)\mathbb{E}[Z^{2} | a_i]}{\sum_{a \in \mathcal{A}} p(a)\mathbb{E}[Z^{2} | a]} \}_{a_i \in \mathcal{A}}$, which can then be used for explorative post-selection by the agent.

\section{Proof of upper bound for cascaded amplitude kickback}

Lower part of upper bound range is $1/|\mathcal{A}|$ (smallest the max value can be). Upper part of range will be some equation related to $n$, and as $n$ tends to infinity the upper part will approach $1$. Upper part will be highly dependant on the initial distribution (wider initial spread will result in lower $n$, narrow spread will need higher $n$).

TODO: How many iterations $n$ until the upper bound approaches $1$?

% https://math.stackexchange.com/questions/386360/upper-bound-on-a-probability-generating-function-with-a-finite-first-moment

% Cramér-Rao lower bound (CRLB) -- https://en.wikipedia.org/wiki/Cramér–Rao_bound
% Moment-based bounds on conditional expectations

\end{document}


% References:

% See:
% https://pennylane.ai/qml/demos/tutorial_mcm_introduction#qubit-reset
% https://quantumcomputing.stackexchange.com/questions/4834/dens/ity-matrix-after-measurement-on-density-matrix
% https://physics.stackexchange.com/questions/681556/trying-to-understand-post-measurement-density-matrices-in-a-state-that-spans-2-h
